\section{Motivation}

The density fields generated by Direct Numerical Simulations of the Rayleigh--Taylor instability contain a large amount of spatial information distributed across multiple scales. As discussed in the previous chapter, the instability evolves from relatively smooth interface perturbations toward highly complex turbulent structures characterized by sharp gradients, coherent vortices, and intermittent mixing regions.

From a machine learning perspective, directly learning probability distributions in such high dimensional spaces is particularly challenging. The dimensionality of the data increases rapidly with the spatial resolution of the simulations, leading to high memory requirements and expensive training procedures. Moreover, the limited number of DNS samples available further complicates the learning process, since the model may tend to memorize the training examples rather than learn patterns that generalize to unseen samples.

Reducing the dimensional complexity of the data while preserving the physically relevant structures therefore constitutes an important objective. Beyond computational considerations, appropriate representations may also improve the stability and efficiency of generative models by emphasizing dominant structures and reducing redundant information.

Another important aspect concerns the strongly multi-scale nature of RTI flows. Large coherent structures coexist with fine turbulent features and localized interface deformations. Consequently, representations capable of simultaneously capturing both large-scale organization and localized high-frequency structures are particularly attractive.

Several approaches can be considered for this purpose. Classical spectral decompositions based on Fourier transforms provide efficient global frequency representations and are widely used in turbulence analysis. However, they remain limited when dealing with localized or strongly discontinuous structures. Multi-resolution methods based on wavelet transforms constitute an alternative approach capable of jointly representing spatial localization and scale-dependent information.

The objective of this chapter is therefore to introduce the different representations investigated in this work and to analyze their respective advantages and limitations in the context of Rayleigh--Taylor instability fields.

